% =============================================================================
% 08-cpu-scheduling-basics.tex — Lecture 8:
% CPU Scheduling: Concepts, Criteria, FCFS and SJF
%
% Course: Operating Systems (PCC CS-502)
% Anchor reading: Silberschatz, Galvin & Gagne, Operating System Concepts,
%                 8th ed. — Ch. 5 §5.1 (basic concepts, p. 183), §5.2
%                 (scheduling criteria, p. 187), §5.3.1 FCFS (p. 188) and
%                 §5.3.2 SJF/SRTF (p. 189); page-verified via
%                 master_supporting_docs/CS502/supporting_books/SilberschatzOSC/index.md.
%                 Stallings, Operating Systems: Internals and Design
%                 Principles, 6th ed. — Ch. 9 §9.1 (types of scheduling,
%                 p. 406) and §9.2 (scheduling algorithms, p. 410);
%                 page-verified via .../StallingsOS/index.md.
%
% Every Gantt chart and every CT/TAT/WT/RT table in this deck was generated
% from a verified simulator (FCFS, non-preemptive SJF, SRTF; ties broken by
% arrival order; SRTF preempts only on a strictly shorter remaining time).
% None of the worked instances contains a tie.
%
% Compile with: /compile-latex CS502/08-cpu-scheduling-basics
% =============================================================================

\documentclass[aspectratio=169]{beamer}
\input{header}
\coursecode{CS502}

\title{CPU Scheduling}
\subtitle{Lecture 8 --- Concepts, Criteria, FCFS and SJF}
\author[Dr. Auqib Hamid Lone]{%
  \textbf{Dr. Auqib Hamid Lone}\\
  \small Assistant Professor in Computer Science \& Engineering
}
\institute[GCET Ganderbal]{%
  \small\textbf{PCC CS-502: Operating Systems}\\[0.3em]
  Department of Computer Science \& Engineering\\
  Government College of Engineering \& Technology (GCET), Ganderbal
}
\date[\today]{\small\today}

\begin{document}

\frame{\titlepage}

% =============================================================================
% ACT 0 --- FIVE JOBS, ONE CPU
% =============================================================================
\section{Five Jobs, One CPU}

\begin{frame}{Welcome: Who Runs Next on the Lab Server?}
  \begin{block}{The big question}
    Five students submit compile jobs to the department's lab server within
    five milliseconds. The server has one CPU. In what order should the jobs
    run --- and how would you even decide that one order is better than
    another?
  \end{block}
  \centering
  \footnotesize
  \begin{tabular}{lccccc}
    \toprule
    Process & \texttt{P1} & \texttt{P2} & \texttt{P3} & \texttt{P4} & \texttt{P5} \\
    \midrule
    Arrival time (ms) & 0 & 1 & 2 & 3 & 5 \\
    CPU burst $\tau$ (ms) & 6 & 8 & 2 & 5 & 4 \\
    \bottomrule
  \end{tabular}
  \vspace{0.2cm}
  \begin{alertblock}{Running thread for today}
    We schedule these five jobs three ways --- FCFS, SJF, SRTF --- and measure
    each schedule with the same yardsticks.
  \end{alertblock}
\end{frame}

\begin{frame}{Roadmap: Four Questions About the Ready Queue}
  \begin{enumerate}
    \item \key{Why schedule at all?} --- the CPU--I/O burst cycle, the
      scheduler, the dispatcher, and preemption.
    \item \key{What makes a schedule good?} --- five criteria and the formulas
      for TAT, WT, and RT.
    \item \key{What does the simplest policy do?} --- FCFS and the convoy
      effect.
    \item \key{Can we do better?} --- SJF, its optimality, burst prediction,
      and the preemptive form SRTF.
  \end{enumerate}
  \vspace{0.2cm}
  \muted{Lecture 6 introduced the short-term scheduler as a box in the state
  diagram. Today we open the box and ask what policy it should run.}
\end{frame}

\transitionslide{A process does not use the CPU continuously. It alternates between computing and waiting --- and the waiting is what makes scheduling worthwhile.}

% =============================================================================
% ACT 1 --- BASIC CONCEPTS
% =============================================================================
\section{Basic Concepts}

\begin{frame}{The CPU--I/O Burst Cycle}
  \centering
  \begin{tikzpicture}[every node/.style={font=\small}]
    % Coordinate map: six boxes on y=0, centres x = -5.5, -3.3, -1.1, 1.1, 3.3, 5.5;
    % each box 1.9cm wide (gap 0.3cm), extent x = -6.45..6.45. CPU boxes use
    % dag-node, I/O boxes dag-node with gold border, white fill. Time arrow on y=-0.9, x = -6.45..6.45.
    \node[dag-node, minimum width=1.9cm, minimum height=0.8cm] (c1) at (-5.5,0) {CPU burst};
    \node[dag-node, draw=primary-gold, fill=white, minimum width=1.9cm, minimum height=0.8cm] (i1) at (-3.3,0) {I/O burst};
    \node[dag-node, minimum width=1.9cm, minimum height=0.8cm] (c2) at (-1.1,0) {CPU burst};
    \node[dag-node, draw=primary-gold, fill=white, minimum width=1.9cm, minimum height=0.8cm] (i2) at (1.1,0) {I/O burst};
    \node[dag-node, minimum width=1.9cm, minimum height=0.8cm] (c3) at (3.3,0) {CPU burst};
    \node[dag-node, draw=primary-gold, fill=white, minimum width=1.9cm, minimum height=0.8cm] (i3) at (5.5,0) {\ldots};
    \draw[observed-edge] (-6.45,-0.9) -- node[below, font=\scriptsize] {time} (6.45,-0.9);
  \end{tikzpicture}
  \vspace{0.25cm}
  \begin{itemize}
    \item Execution is a cycle: a \key{CPU burst}, then an \key{I/O burst},
      then another CPU burst, until a final CPU burst ends in termination.
    \item Measured burst lengths are dominated by many short bursts and a
      few long ones.
    \item An \key{I/O-bound} program has many short CPU bursts; a
      \key{CPU-bound} program has a few long ones.
  \end{itemize}
  \muted{Silberschatz et al., Ch.~5 \S5.1, p.~183.}
\end{frame}

\begin{frame}{The CPU Scheduler and the Ready Queue}
  \begin{block}{Definition}
    The \key{CPU scheduler} (short-term scheduler) selects one process from
    the processes in memory that are ready to execute and allocates the CPU
    to it.
  \end{block}
  \begin{itemize}
    \item The ready queue holds PCBs, not programs. It is \emph{not}
      necessarily FIFO: it can be a FIFO queue, a priority queue, a tree, or
      an unordered list.
    \item The \emph{policy} is the rule that picks from it. FCFS, SJF, and
      SRTF are three such rules.
    \item While one process waits for I/O, the scheduler gives the CPU to
      another --- that is how multiprogramming keeps the CPU busy.
  \end{itemize}
  \muted{The long-term and medium-term schedulers of Lecture 6 decide
  \emph{who is in memory}; today's scheduler decides \emph{who runs}.}
\end{frame}

\begin{frame}{When Does the Scheduler Run? Preemptive vs.\ Non-preemptive}
  \footnotesize
  A scheduling decision may be needed when a process:
  \begin{enumerate}
    \item switches from running to waiting (an I/O request, \texttt{wait});
    \item switches from running to ready (a timer interrupt);
    \item switches from waiting to ready (I/O completion);
    \item terminates.
  \end{enumerate}
  \begin{columns}[T]
    \begin{column}{0.48\textwidth}
      \begin{block}{Non-preemptive (cooperative)}
        Decisions only in cases 1 and 4. Once a process has the CPU, it keeps
        it until it blocks or exits.
      \end{block}
    \end{column}
    \begin{column}{0.48\textwidth}
      \begin{alertblock}{Preemptive}
        Decisions also in cases 2 and 3. The kernel may take the CPU away ---
        and must then protect shared data from a half-finished update.
      \end{alertblock}
    \end{column}
  \end{columns}
  \muted{Silberschatz et al., Ch.~5 \S5.1, pp.~183--187.}
\end{frame}

\begin{frame}{The Dispatcher and Dispatch Latency}
  \begin{block}{Definition}
    The \key{dispatcher} is the module that gives control of the CPU to the
    process the scheduler selected. The time it takes to stop one process and
    start another is the \key{dispatch latency}.
  \end{block}
  \begin{itemize}
    \item It switches context (saves one PCB, loads another --- Lecture 6).
    \item It switches to user mode.
    \item It jumps to the saved PC in the new process.
  \end{itemize}
  \vspace{0.1cm}
  \begin{exampleblock}{Why it matters}
    The dispatcher runs on \emph{every} switch. A policy that switches more
    often pays the dispatch latency more often --- we will measure that cost
    on the lab-server jobs.
  \end{exampleblock}
\end{frame}

\begin{frame}{Socratic Check 1}
  \begin{block}{Question}
    A process finishes its CPU burst and issues a disk \texttt{read}. Under a
    \emph{non-preemptive} scheduler, must a scheduling decision be made? What
    about when that \texttt{read} later completes?
  \end{block}
  \vspace{0.25cm}
  \begin{itemize}
    \item Issuing the \texttt{read} is case 1 (running $\to$ waiting): the CPU
      is free, so the scheduler must pick someone.
    \item The completion is case 3 (waiting $\to$ ready): a non-preemptive
      scheduler lets the current process keep the CPU.
    \item Only a preemptive scheduler may reconsider at case 3.
  \end{itemize}
\end{frame}

\transitionslide{Before comparing policies we need yardsticks. ``Better'' means nothing until we say better for whom, measured how.}

% =============================================================================
% ACT 2 --- SCHEDULING CRITERIA
% =============================================================================
\section{Scheduling Criteria}

\begin{frame}{Five Criteria}
  \centering
  \footnotesize
  \renewcommand{\arraystretch}{1.3}
  \begin{tabular}{p{0.22\textwidth}p{0.50\textwidth}p{0.16\textwidth}}
    \toprule
    \textbf{Criterion} & \textbf{Meaning} & \textbf{Goal} \\
    \midrule
    CPU utilization & fraction of time the CPU does useful work & maximize \\
    Throughput & processes completed per unit time & maximize \\
    Turnaround time (TAT) & submission to completion, per process & minimize \\
    Waiting time (WT) & total time spent in the ready queue & minimize \\
    Response time (RT) & submission to the \emph{first} time on the CPU & minimize \\
    \bottomrule
  \end{tabular}
  \vspace{0.25cm}
  \begin{itemize}
    \item The scheduling policy directly controls only \key{waiting time}; it
      does not change how long a burst computes.
    \item We usually optimize the \emph{average}; interactive systems also
      care about the \emph{variance} of response time.
  \end{itemize}
  \muted{Silberschatz et al., Ch.~5 \S5.2, pp.~187--188.}
\end{frame}

\begin{frame}{The Three Per-Process Formulas}
  \begin{block}{Definitions (one CPU burst per process)}
    For process \texttt{Pi} with arrival time $\mathrm{AT}_i$, burst
    $\tau_i$, first-dispatch time $\mathrm{ST}_i$, and completion time
    $\mathrm{CT}_i$:
    \[
      \mathrm{TAT}_i = \mathrm{CT}_i - \mathrm{AT}_i, \qquad
      \mathrm{WT}_i = \mathrm{TAT}_i - \tau_i, \qquad
      \mathrm{RT}_i = \mathrm{ST}_i - \mathrm{AT}_i .
    \]
  \end{block}
  \begin{itemize}
    \item WT counts \emph{all} time ready but not running --- including time
      after a preemption.
    \item For a non-preemptive policy, the process runs once, so
      $\mathrm{WT}_i = \mathrm{RT}_i$.
    \item Idle CPU time is not anybody's waiting time: no process is ready.
  \end{itemize}
  \muted{TAT, WT, and RT are spelled out here and abbreviated from now on.}
\end{frame}

\begin{frame}{Worked Example: Metrics From a Schedule With an Idle Gap}
  \centering
  \footnotesize
  Jobs: \texttt{P1} $(0, 3)$, \texttt{P2} $(5, 4)$, \texttt{P3} $(6, 2)$ as
  (AT, $\tau$) in ms, scheduled FCFS.\\[0.2cm]
  \begin{tikzpicture}[x=0.75cm,y=1cm]
  % Gantt: I under FCFS (generated from the verified simulator); time axis left to right,
  % bars y=0..0.6, time labels below y=0, arrival arrows y=-0.95..-0.5, arrival names below y=-0.95
  \draw[draw=primary-blue, fill=light-bg] (0,0) rectangle (3,0.6);
  \node[font=\scriptsize] at (1.5,0.3) {\texttt{P1}};
  \draw[draw=neutral, fill=neutral!12] (3,0) rectangle (5,0.6);
  \node[font=\tiny, text=neutral] at (4,0.3) {idle};
  \draw[draw=primary-blue, fill=light-bg] (5,0) rectangle (9,0.6);
  \node[font=\scriptsize] at (7,0.3) {\texttt{P2}};
  \draw[draw=primary-blue, fill=light-bg] (9,0) rectangle (11,0.6);
  \node[font=\scriptsize] at (10,0.3) {\texttt{P3}};
  \node[font=\tiny, below] at (0,0) {0};
  \node[font=\tiny, below] at (3,0) {3};
  \node[font=\tiny, below] at (5,0) {5};
  \node[font=\tiny, below] at (9,0) {9};
  \node[font=\tiny, below] at (11,0) {11};
  \draw[-{Latex[length=1.2mm]}, draw=primary-gold] (0,-0.95) -- (0,-0.5);
  \node[font=\tiny, text=primary-gold, below] at (0,-0.95) {P1};
  \draw[-{Latex[length=1.2mm]}, draw=primary-gold] (5,-0.95) -- (5,-0.5);
  \node[font=\tiny, text=primary-gold, below] at (5,-0.95) {P2};
  \draw[-{Latex[length=1.2mm]}, draw=primary-gold] (6,-0.95) -- (6,-0.5);
  \node[font=\tiny, text=primary-gold, below] at (6,-0.95) {P3};
\end{tikzpicture}
  \\[0.15cm]
  {\footnotesize\begin{tabular}{lrrrrrr}
\toprule
\textbf{Proc.} & \textbf{AT} & $\boldsymbol{\tau}$ & \textbf{CT} & \textbf{TAT} & \textbf{WT} & \textbf{RT} \\
\midrule
\texttt{P1} & 0 & 3 & 3 & 3 & 0 & 0 \\
\texttt{P2} & 5 & 4 & 9 & 4 & 0 & 0 \\
\texttt{P3} & 6 & 2 & 11 & 5 & 3 & 3 \\
\midrule
\multicolumn{4}{@{}l}{\textbf{Average}} & 4 & 1 & 1 \\
\bottomrule
\end{tabular}}
  \vspace{0.15cm}
  \begin{itemize}
    \item The CPU idles from 3 to 5: nothing has arrived. \texttt{P2} starts
      at its arrival, so $\mathrm{WT}_2 = 0$.
    \item \texttt{P3}: $\mathrm{TAT}_3 = 11 - 6 = 5$, $\mathrm{WT}_3 = 5 - 2 = 3$.
    \item Utilization over $[0, 11]$ is $9/11 \approx 81.8\%$.
  \end{itemize}
\end{frame}

\begin{frame}{Utilization and Throughput Include the Dispatcher}
  \begin{itemize}
    \item Lab-server jobs: total burst $6+8+2+5+4 = 25$~ms.
    \item FCFS dispatches each job once: 5 dispatches.
    \item With a dispatch latency of $0.5$~ms per dispatch, the CPU spends
      $25 + 5 \times 0.5 = 27.5$~ms to finish all five.
  \end{itemize}
  \vspace{0.15cm}
  \begin{exampleblock}{Worked numbers}
    \[
      \text{utilization} = \frac{25}{27.5} \approx 90.9\%, \qquad
      \text{throughput} = \frac{5\ \text{jobs}}{27.5\ \text{ms}} \approx 0.18\ \text{jobs/ms}.
    \]
  \end{exampleblock}
  \muted{With zero dispatch latency the same schedule gives $100\%$ and
  $5/25 = 0.2$ jobs/ms. Every extra switch moves both numbers down.}
\end{frame}

\begin{frame}{Socratic Check 2}
  \begin{block}{Question}
    For a laptop running a text editor, a browser, and a long video export,
    which criterion would you optimize first --- and why is ``maximize
    throughput'' the wrong answer?
  \end{block}
  \vspace{0.25cm}
  \begin{itemize}
    \item The user feels \key{response time}: the editor must react to a
      keystroke within milliseconds.
    \item Maximizing throughput could run the export to completion first and
      leave the editor frozen.
    \item Criteria conflict; a policy is a choice of which one to favour.
  \end{itemize}
\end{frame}

\transitionslide{The simplest rule is the one at every ticket counter: serve whoever arrived first. It is fair-sounding, easy to build --- and can be very slow.}

% =============================================================================
% ACT 3 --- FCFS
% =============================================================================
\section{First-Come, First-Served}

\begin{frame}{FCFS: Definition}
  \begin{block}{Definition}
    \key{First-Come, First-Served} (FCFS) allocates the CPU to the process
    that requested it first. The ready queue is a FIFO queue: an arriving
    PCB joins the tail; the scheduler takes the head.
  \end{block}
  \begin{itemize}
    \item FCFS is \key{non-preemptive}: the running process keeps the CPU
      until it blocks or terminates.
    \item The code is a linked list with a head and a tail pointer.
    \item Its average waiting time is often long --- and depends heavily on
      the order of arrival.
  \end{itemize}
  \muted{Silberschatz et al., Ch.~5 \S5.3.1, pp.~188--189; Stallings,
  Ch.~9 \S9.2, p.~410.}
\end{frame}

\begin{frame}{Worked Example: One Long Job First}
  \centering
  \footnotesize
  \texttt{P1}, \texttt{P2}, \texttt{P3} arrive at time 0 in that order, with
  bursts 24, 3, 3~ms.\\[0.25cm]
  \begin{tikzpicture}[x=0.4cm,y=1cm]
  % Gantt: S1 under FCFS (generated from the verified simulator); time axis left to right,
  % bars y=0..0.6, time labels below y=0, arrival arrows y=-0.95..-0.5, arrival names below y=-0.95
  \draw[draw=primary-blue, fill=light-bg] (0,0) rectangle (24,0.6);
  \node[font=\scriptsize] at (12,0.3) {\texttt{P1}};
  \draw[draw=primary-blue, fill=light-bg] (24,0) rectangle (27,0.6);
  \node[font=\scriptsize] at (25.5,0.3) {\texttt{P2}};
  \draw[draw=primary-blue, fill=light-bg] (27,0) rectangle (30,0.6);
  \node[font=\scriptsize] at (28.5,0.3) {\texttt{P3}};
  \node[font=\tiny, below] at (0,0) {0};
  \node[font=\tiny, below] at (24,0) {24};
  \node[font=\tiny, below] at (27,0) {27};
  \node[font=\tiny, below] at (30,0) {30};
\end{tikzpicture}
  \\[0.25cm]
  \begin{tabular}{lccc}
    \toprule
    & \texttt{P1} & \texttt{P2} & \texttt{P3} \\
    \midrule
    WT (ms) & 0 & 24 & 27 \\
    \bottomrule
  \end{tabular}
  \vspace{0.2cm}
  \begin{alertblock}{Average WT}
    $(0 + 24 + 27)/3 = 17$~ms. Two 3-ms jobs wait behind one 24-ms job.
  \end{alertblock}
\end{frame}

\begin{frame}{Same Jobs, Different Arrival Order}
  \centering
  \footnotesize
  Now \texttt{P2}, \texttt{P3}, \texttt{P1} arrive at time 0 in that order.\\[0.25cm]
  \begin{tikzpicture}[x=0.4cm,y=1cm]
  % Gantt: S1b under FCFS (generated from the verified simulator); time axis left to right,
  % bars y=0..0.6, time labels below y=0, arrival arrows y=-0.95..-0.5, arrival names below y=-0.95
  \draw[draw=primary-blue, fill=light-bg] (0,0) rectangle (3,0.6);
  \node[font=\scriptsize] at (1.5,0.3) {\texttt{P2}};
  \draw[draw=primary-blue, fill=light-bg] (3,0) rectangle (6,0.6);
  \node[font=\scriptsize] at (4.5,0.3) {\texttt{P3}};
  \draw[draw=primary-blue, fill=light-bg] (6,0) rectangle (30,0.6);
  \node[font=\scriptsize] at (18,0.3) {\texttt{P1}};
  \node[font=\tiny, below] at (0,0) {0};
  \node[font=\tiny, below] at (3,0) {3};
  \node[font=\tiny, below] at (6,0) {6};
  \node[font=\tiny, below] at (30,0) {30};
\end{tikzpicture}
  \\[0.25cm]
  \begin{tabular}{lccc}
    \toprule
    & \texttt{P1} & \texttt{P2} & \texttt{P3} \\
    \midrule
    WT (ms) & 6 & 0 & 3 \\
    \bottomrule
  \end{tabular}
  \vspace{0.2cm}
  \begin{exampleblock}{Average WT}
    $(6 + 0 + 3)/3 = 3$~ms --- a factor of $17/3 \approx 5.7$ better, with
    the same work and the same finish time of 30~ms.
  \end{exampleblock}
\end{frame}

\begin{frame}{The Convoy Effect}
  \begin{itemize}
    \item One \key{CPU-bound} process holds the CPU for a long burst.
    \item Many \key{I/O-bound} processes finish their I/O and queue behind
      it; the I/O devices sit idle.
    \item The CPU-bound process finally blocks; the short jobs run quickly
      and block on I/O; now the CPU sits idle.
    \item The cycle repeats: everyone moves at the speed of the slowest.
  \end{itemize}
  \vspace{0.2cm}
  \begin{alertblock}{Consequence}
    FCFS gives lower CPU and device utilization than letting the short jobs
    go first --- and, being non-preemptive, it is unsuitable for time-sharing
    systems.
  \end{alertblock}
  \muted{Silberschatz et al., Ch.~5 \S5.3.1, pp.~188--189.}
\end{frame}

\begin{frame}{The Lab Server Under FCFS}
  \centering
  \begin{tikzpicture}[x=0.5cm,y=1cm]
  % Gantt: R under FCFS (generated from the verified simulator); time axis left to right,
  % bars y=0..0.6, time labels below y=0, arrival arrows y=-0.95..-0.5, arrival names below y=-0.95
  \draw[draw=primary-blue, fill=light-bg] (0,0) rectangle (6,0.6);
  \node[font=\scriptsize] at (3,0.3) {\texttt{P1}};
  \draw[draw=primary-blue, fill=light-bg] (6,0) rectangle (14,0.6);
  \node[font=\scriptsize] at (10,0.3) {\texttt{P2}};
  \draw[draw=primary-blue, fill=light-bg] (14,0) rectangle (16,0.6);
  \node[font=\scriptsize] at (15,0.3) {\texttt{P3}};
  \draw[draw=primary-blue, fill=light-bg] (16,0) rectangle (21,0.6);
  \node[font=\scriptsize] at (18.5,0.3) {\texttt{P4}};
  \draw[draw=primary-blue, fill=light-bg] (21,0) rectangle (25,0.6);
  \node[font=\scriptsize] at (23,0.3) {\texttt{P5}};
  \node[font=\tiny, below] at (0,0) {0};
  \node[font=\tiny, below] at (6,0) {6};
  \node[font=\tiny, below] at (14,0) {14};
  \node[font=\tiny, below] at (16,0) {16};
  \node[font=\tiny, below] at (21,0) {21};
  \node[font=\tiny, below] at (25,0) {25};
  \draw[-{Latex[length=1.2mm]}, draw=primary-gold] (0,-0.95) -- (0,-0.5);
  \node[font=\tiny, text=primary-gold, below] at (0,-0.95) {P1};
  \draw[-{Latex[length=1.2mm]}, draw=primary-gold] (1,-0.95) -- (1,-0.5);
  \node[font=\tiny, text=primary-gold, below] at (1,-0.95) {P2};
  \draw[-{Latex[length=1.2mm]}, draw=primary-gold] (2,-0.95) -- (2,-0.5);
  \node[font=\tiny, text=primary-gold, below] at (2,-0.95) {P3};
  \draw[-{Latex[length=1.2mm]}, draw=primary-gold] (3,-0.95) -- (3,-0.5);
  \node[font=\tiny, text=primary-gold, below] at (3,-0.95) {P4};
  \draw[-{Latex[length=1.2mm]}, draw=primary-gold] (5,-0.95) -- (5,-0.5);
  \node[font=\tiny, text=primary-gold, below] at (5,-0.95) {P5};
\end{tikzpicture}
  \\[0.2cm]
  {\footnotesize\begin{tabular}{lrrrrrr}
\toprule
\textbf{Proc.} & \textbf{AT} & $\boldsymbol{\tau}$ & \textbf{CT} & \textbf{TAT} & \textbf{WT} & \textbf{RT} \\
\midrule
\texttt{P1} & 0 & 6 & 6 & 6 & 0 & 0 \\
\texttt{P2} & 1 & 8 & 14 & 13 & 5 & 5 \\
\texttt{P3} & 2 & 2 & 16 & 14 & 12 & 12 \\
\texttt{P4} & 3 & 5 & 21 & 18 & 13 & 13 \\
\texttt{P5} & 5 & 4 & 25 & 20 & 16 & 16 \\
\midrule
\multicolumn{4}{@{}l}{\textbf{Average}} & 14.2 & 9.2 & 9.2 \\
\bottomrule
\end{tabular}}
  \vspace{0.15cm}
  \begin{itemize}
    \footnotesize
    \item \texttt{P3} needs only 2~ms but waits 12~ms behind \texttt{P2}'s
      8-ms burst. Average WT $= 46/5 = 9.2$~ms.
  \end{itemize}
\end{frame}

\transitionslide{If order matters so much, pick the order deliberately: run the shortest job first.}

% =============================================================================
% ACT 4 --- SJF
% =============================================================================
\section{Shortest-Job-First}

\begin{frame}{SJF: Definition}
  \begin{block}{Definition}
    \key{Shortest-Job-First} (SJF) gives the CPU to the ready process with
    the smallest \emph{next CPU burst}. Ties are broken FCFS.
  \end{block}
  \begin{itemize}
    \item The more accurate name is \emph{shortest-next-CPU-burst}: the
      length that matters is the next burst, not the whole job.
    \item In the form we use first, SJF is \key{non-preemptive}: once chosen,
      a process runs its whole burst.
    \item The preemptive form is SRTF (later today).
  \end{itemize}
  \muted{Silberschatz et al., Ch.~5 \S5.3.2, pp.~189--192; Stallings calls
  it SPN (shortest process next), Ch.~9 \S9.2.}
\end{frame}

\begin{frame}{Worked Example: SJF With Simultaneous Arrivals}
  \centering
  \footnotesize
  \texttt{P1}--\texttt{P4} arrive at 0 with bursts 6, 8, 7, 3~ms.\\[0.25cm]
  \begin{tikzpicture}[x=0.5cm,y=1cm]
  % Gantt: S2 under SJF (generated from the verified simulator); time axis left to right,
  % bars y=0..0.6, time labels below y=0, arrival arrows y=-0.95..-0.5, arrival names below y=-0.95
  \draw[draw=primary-blue, fill=light-bg] (0,0) rectangle (3,0.6);
  \node[font=\scriptsize] at (1.5,0.3) {\texttt{P4}};
  \draw[draw=primary-blue, fill=light-bg] (3,0) rectangle (9,0.6);
  \node[font=\scriptsize] at (6,0.3) {\texttt{P1}};
  \draw[draw=primary-blue, fill=light-bg] (9,0) rectangle (16,0.6);
  \node[font=\scriptsize] at (12.5,0.3) {\texttt{P3}};
  \draw[draw=primary-blue, fill=light-bg] (16,0) rectangle (24,0.6);
  \node[font=\scriptsize] at (20,0.3) {\texttt{P2}};
  \node[font=\tiny, below] at (0,0) {0};
  \node[font=\tiny, below] at (3,0) {3};
  \node[font=\tiny, below] at (9,0) {9};
  \node[font=\tiny, below] at (16,0) {16};
  \node[font=\tiny, below] at (24,0) {24};
\end{tikzpicture}
  \\[0.25cm]
  \begin{tabular}{lcccc}
    \toprule
    & \texttt{P1} & \texttt{P2} & \texttt{P3} & \texttt{P4} \\
    \midrule
    WT (ms) & 3 & 16 & 9 & 0 \\
    \bottomrule
  \end{tabular}
  \vspace{0.2cm}
  \begin{exampleblock}{Comparison}
    SJF: $(3+16+9+0)/4 = 7$~ms. FCFS in index order:
    $(0+6+14+21)/4 = 10.25$~ms.
  \end{exampleblock}
\end{frame}

\begin{frame}{Why SJF Is Optimal}
  \begin{itemize}
    \item $n$ processes all ready at time 0, run in order with bursts
      $\tau_{(1)}, \tau_{(2)}, \ldots, \tau_{(n)}$.
    \item The $k$-th burst delays every one of the $n-k$ processes after it:
      \[
        \sum_{i} \mathrm{WT}_i = \sum_{k=1}^{n} (n-k)\,\tau_{(k)} .
      \]
    \item The largest weight $n-1$ multiplies the first burst. To minimize
      the sum, put the smallest burst there --- and so on down.
  \end{itemize}
  \begin{alertblock}{Result}
    For a given set of processes, SJF gives the minimum average waiting
    time. Swapping any adjacent long--short pair into short--long order
    lowers the total.
  \end{alertblock}
  \muted{Silberschatz et al., Ch.~5 \S5.3.2, from p.~189 (SJF is provably optimal).}
\end{frame}

\begin{frame}{The Lab Server Under Non-preemptive SJF}
  \centering
  \begin{tikzpicture}[x=0.5cm,y=1cm]
  % Gantt: R under SJF (generated from the verified simulator); time axis left to right,
  % bars y=0..0.6, time labels below y=0, arrival arrows y=-0.95..-0.5, arrival names below y=-0.95
  \draw[draw=primary-blue, fill=light-bg] (0,0) rectangle (6,0.6);
  \node[font=\scriptsize] at (3,0.3) {\texttt{P1}};
  \draw[draw=primary-blue, fill=light-bg] (6,0) rectangle (8,0.6);
  \node[font=\scriptsize] at (7,0.3) {\texttt{P3}};
  \draw[draw=primary-blue, fill=light-bg] (8,0) rectangle (12,0.6);
  \node[font=\scriptsize] at (10,0.3) {\texttt{P5}};
  \draw[draw=primary-blue, fill=light-bg] (12,0) rectangle (17,0.6);
  \node[font=\scriptsize] at (14.5,0.3) {\texttt{P4}};
  \draw[draw=primary-blue, fill=light-bg] (17,0) rectangle (25,0.6);
  \node[font=\scriptsize] at (21,0.3) {\texttt{P2}};
  \node[font=\tiny, below] at (0,0) {0};
  \node[font=\tiny, below] at (6,0) {6};
  \node[font=\tiny, below] at (8,0) {8};
  \node[font=\tiny, below] at (12,0) {12};
  \node[font=\tiny, below] at (17,0) {17};
  \node[font=\tiny, below] at (25,0) {25};
  \draw[-{Latex[length=1.2mm]}, draw=primary-gold] (0,-0.95) -- (0,-0.5);
  \node[font=\tiny, text=primary-gold, below] at (0,-0.95) {P1};
  \draw[-{Latex[length=1.2mm]}, draw=primary-gold] (1,-0.95) -- (1,-0.5);
  \node[font=\tiny, text=primary-gold, below] at (1,-0.95) {P2};
  \draw[-{Latex[length=1.2mm]}, draw=primary-gold] (2,-0.95) -- (2,-0.5);
  \node[font=\tiny, text=primary-gold, below] at (2,-0.95) {P3};
  \draw[-{Latex[length=1.2mm]}, draw=primary-gold] (3,-0.95) -- (3,-0.5);
  \node[font=\tiny, text=primary-gold, below] at (3,-0.95) {P4};
  \draw[-{Latex[length=1.2mm]}, draw=primary-gold] (5,-0.95) -- (5,-0.5);
  \node[font=\tiny, text=primary-gold, below] at (5,-0.95) {P5};
\end{tikzpicture}
  \\[0.2cm]
  {\footnotesize\begin{tabular}{lrrrrrr}
\toprule
\textbf{Proc.} & \textbf{AT} & $\boldsymbol{\tau}$ & \textbf{CT} & \textbf{TAT} & \textbf{WT} & \textbf{RT} \\
\midrule
\texttt{P1} & 0 & 6 & 6 & 6 & 0 & 0 \\
\texttt{P2} & 1 & 8 & 25 & 24 & 16 & 16 \\
\texttt{P3} & 2 & 2 & 8 & 6 & 4 & 4 \\
\texttt{P4} & 3 & 5 & 17 & 14 & 9 & 9 \\
\texttt{P5} & 5 & 4 & 12 & 7 & 3 & 3 \\
\midrule
\multicolumn{4}{@{}l}{\textbf{Average}} & 11.4 & 6.4 & 6.4 \\
\bottomrule
\end{tabular}}
  \vspace{0.15cm}
  \begin{itemize}
    \footnotesize
    \item At $t=0$ only \texttt{P1} is ready, so it runs to 6. At 6 the
      choice is $\{8, 2, 5, 4\}$: \texttt{P3}, then \texttt{P5}, \texttt{P4},
      \texttt{P2}. Average WT $= 32/5 = 6.4$~ms.
  \end{itemize}
\end{frame}

\begin{frame}{The Catch: Nobody Knows the Next Burst}
  \begin{itemize}
    \item The short-term scheduler cannot read the future. It can only
      \emph{predict} the next CPU burst from the past ones.
    \item The standard predictor is the \key{exponential average}. With
      $\tau_n$ the measured length of the $n$-th burst and
      $\hat{\tau}_n$ its prediction:
  \end{itemize}
  \begin{block}{Exponential averaging}
    \[
      \hat{\tau}_{n+1} = \alpha\,\tau_n + (1-\alpha)\,\hat{\tau}_n ,
      \qquad 0 \le \alpha \le 1 .
    \]
  \end{block}
  \begin{itemize}
    \item $\alpha = 0$: history only, recent behaviour ignored.
      $\alpha = 1$: only the last burst counts.
    \item Expanding the recursion, older bursts get weights $\alpha(1-\alpha)^j$
      --- each step back counts less.
  \end{itemize}
  \muted{Silberschatz writes $t_n$ for the measured burst; we keep $\tau$ for
  bursts and add a hat for a prediction.}
\end{frame}

\begin{frame}{Worked Example: Predicting Bursts With $\alpha = 1/2$}
  \centering
  \footnotesize
  $\hat{\tau}_0 = 10$~ms; measured bursts $6, 4, 6, 4, 13, 13, 13$~ms.\\[0.25cm]
  \renewcommand{\arraystretch}{1.25}
  \begin{tabular}{lcccccccc}
    \toprule
    $n$ & 0 & 1 & 2 & 3 & 4 & 5 & 6 & 7 \\
    \midrule
    measured $\tau_n$ & 6 & 4 & 6 & 4 & 13 & 13 & 13 & --- \\
    predicted $\hat{\tau}_n$ & 10 & 8 & 6 & 6 & 5 & 9 & 11 & 12 \\
    \bottomrule
  \end{tabular}
  \vspace{0.25cm}
  \begin{itemize}
    \item $\hat{\tau}_1 = \tfrac12(6) + \tfrac12(10) = 8$;
      $\hat{\tau}_2 = \tfrac12(4) + \tfrac12(8) = 6$;
      $\hat{\tau}_5 = \tfrac12(13) + \tfrac12(5) = 9$.
    \item When behaviour jumps from 4 to 13, the prediction lags --- 9, 11,
      12 --- and converges.
  \end{itemize}
  \muted{Same instance as the exponential-average illustration in Silberschatz et al., Ch.~5 \S5.3.2.}
\end{frame}

\begin{frame}{Socratic Check 3}
  \begin{block}{Question}
    A process alternates bursts of 2~ms and 20~ms. Would you choose
    $\alpha$ close to 0 or close to 1, and what does SJF then do with it?
  \end{block}
  \vspace{0.25cm}
  \begin{itemize}
    \item $\alpha$ near 1 follows the last burst: it predicts 2 right before a
      20 and 20 right before a 2 --- always wrong.
    \item $\alpha$ near 0 settles near the long-run average (about 11~ms):
      never exact, but stable.
    \item SJF is only as good as its predictor; a bad guess turns ``shortest
      first'' into an arbitrary order.
  \end{itemize}
\end{frame}

\transitionslide{Non-preemptive SJF ignores a short job that arrives one millisecond after a long one starts. The preemptive form does not.}

% =============================================================================
% ACT 5 --- SRTF
% =============================================================================
\section{Shortest-Remaining-Time-First}

\begin{frame}{SRTF: Preemptive SJF}
  \begin{block}{Definition}
    \key{Shortest-Remaining-Time-First} (SRTF) is preemptive SJF. When a
    process arrives, its burst is compared with the \emph{remaining} time of
    the running process; if it is strictly shorter, the running process is
    preempted.
  \end{block}
  \begin{itemize}
    \item Decisions happen at every arrival and every completion (cases 1,
      3, and 4).
    \item Equal remaining times do not preempt; ties go to the earlier
      arrival.
    \item WT now includes the time a preempted process spends back in the
      ready queue.
  \end{itemize}
  \muted{Silberschatz et al., Ch.~5 \S5.3.2 (from p.~189); Stallings calls
  it SRT, Ch.~9 \S9.2.}
\end{frame}

\begin{frame}{Worked Example: SRTF vs.\ Non-preemptive SJF}
  \centering
  \footnotesize
  (AT, $\tau$): \texttt{P1} $(0,8)$, \texttt{P2} $(1,4)$, \texttt{P3} $(2,9)$,
  \texttt{P4} $(3,5)$.\\[0.15cm]
  \begin{tikzpicture}[x=0.42cm,y=1cm]
  % Gantt: S3 under SRTF (generated from the verified simulator); time axis left to right,
  % bars y=0..0.6, time labels below y=0, arrival arrows y=-0.95..-0.5, arrival names below y=-0.95
  \draw[draw=primary-blue, fill=light-bg] (0,0) rectangle (1,0.6);
  \node[font=\scriptsize] at (0.5,0.3) {\texttt{P1}};
  \draw[draw=primary-blue, fill=light-bg] (1,0) rectangle (5,0.6);
  \node[font=\scriptsize] at (3,0.3) {\texttt{P2}};
  \draw[draw=primary-blue, fill=light-bg] (5,0) rectangle (10,0.6);
  \node[font=\scriptsize] at (7.5,0.3) {\texttt{P4}};
  \draw[draw=primary-blue, fill=light-bg] (10,0) rectangle (17,0.6);
  \node[font=\scriptsize] at (13.5,0.3) {\texttt{P1}};
  \draw[draw=primary-blue, fill=light-bg] (17,0) rectangle (26,0.6);
  \node[font=\scriptsize] at (21.5,0.3) {\texttt{P3}};
  \node[font=\tiny, below] at (0,0) {0};
  \node[font=\tiny, below] at (1,0) {1};
  \node[font=\tiny, below] at (5,0) {5};
  \node[font=\tiny, below] at (10,0) {10};
  \node[font=\tiny, below] at (17,0) {17};
  \node[font=\tiny, below] at (26,0) {26};
  \draw[-{Latex[length=1.2mm]}, draw=primary-gold] (0,-0.95) -- (0,-0.5);
  \node[font=\tiny, text=primary-gold, below] at (0,-0.95) {P1};
  \draw[-{Latex[length=1.2mm]}, draw=primary-gold] (1,-0.95) -- (1,-0.5);
  \node[font=\tiny, text=primary-gold, below] at (1,-0.95) {P2};
  \draw[-{Latex[length=1.2mm]}, draw=primary-gold] (2,-0.95) -- (2,-0.5);
  \node[font=\tiny, text=primary-gold, below] at (2,-0.95) {P3};
  \draw[-{Latex[length=1.2mm]}, draw=primary-gold] (3,-0.95) -- (3,-0.5);
  \node[font=\tiny, text=primary-gold, below] at (3,-0.95) {P4};
\end{tikzpicture}
  \\[0.1cm]
  {\scriptsize\begin{tabular}{lrrrrrr}
\toprule
\textbf{Proc.} & \textbf{AT} & $\boldsymbol{\tau}$ & \textbf{CT} & \textbf{TAT} & \textbf{WT} & \textbf{RT} \\
\midrule
\texttt{P1} & 0 & 8 & 17 & 17 & 9 & 0 \\
\texttt{P2} & 1 & 4 & 5 & 4 & 0 & 0 \\
\texttt{P3} & 2 & 9 & 26 & 24 & 15 & 15 \\
\texttt{P4} & 3 & 5 & 10 & 7 & 2 & 2 \\
\midrule
\multicolumn{4}{@{}l}{\textbf{Average}} & 13 & 6.5 & 4.25 \\
\bottomrule
\end{tabular}}
  \vspace{0.1cm}
  \begin{itemize}
    \item At 1, \texttt{P2}'s 4 beats \texttt{P1}'s remaining 7: preempt.
      \texttt{P1} resumes at 10.
    \item SRTF average WT $= 26/4 = 6.5$~ms; non-preemptive SJF gives
      $7.75$~ms.
  \end{itemize}
\end{frame}

\begin{frame}{The Lab Server Under SRTF}
  \centering
  \begin{tikzpicture}[x=0.5cm,y=1cm]
  % Gantt: R under SRTF (generated from the verified simulator); time axis left to right,
  % bars y=0..0.6, time labels below y=0, arrival arrows y=-0.95..-0.5, arrival names below y=-0.95
  \draw[draw=primary-blue, fill=light-bg] (0,0) rectangle (2,0.6);
  \node[font=\scriptsize] at (1,0.3) {\texttt{P1}};
  \draw[draw=primary-blue, fill=light-bg] (2,0) rectangle (4,0.6);
  \node[font=\scriptsize] at (3,0.3) {\texttt{P3}};
  \draw[draw=primary-blue, fill=light-bg] (4,0) rectangle (8,0.6);
  \node[font=\scriptsize] at (6,0.3) {\texttt{P1}};
  \draw[draw=primary-blue, fill=light-bg] (8,0) rectangle (12,0.6);
  \node[font=\scriptsize] at (10,0.3) {\texttt{P5}};
  \draw[draw=primary-blue, fill=light-bg] (12,0) rectangle (17,0.6);
  \node[font=\scriptsize] at (14.5,0.3) {\texttt{P4}};
  \draw[draw=primary-blue, fill=light-bg] (17,0) rectangle (25,0.6);
  \node[font=\scriptsize] at (21,0.3) {\texttt{P2}};
  \node[font=\tiny, below] at (0,0) {0};
  \node[font=\tiny, below] at (2,0) {2};
  \node[font=\tiny, below] at (4,0) {4};
  \node[font=\tiny, below] at (8,0) {8};
  \node[font=\tiny, below] at (12,0) {12};
  \node[font=\tiny, below] at (17,0) {17};
  \node[font=\tiny, below] at (25,0) {25};
  \draw[-{Latex[length=1.2mm]}, draw=primary-gold] (0,-0.95) -- (0,-0.5);
  \node[font=\tiny, text=primary-gold, below] at (0,-0.95) {P1};
  \draw[-{Latex[length=1.2mm]}, draw=primary-gold] (1,-0.95) -- (1,-0.5);
  \node[font=\tiny, text=primary-gold, below] at (1,-0.95) {P2};
  \draw[-{Latex[length=1.2mm]}, draw=primary-gold] (2,-0.95) -- (2,-0.5);
  \node[font=\tiny, text=primary-gold, below] at (2,-0.95) {P3};
  \draw[-{Latex[length=1.2mm]}, draw=primary-gold] (3,-0.95) -- (3,-0.5);
  \node[font=\tiny, text=primary-gold, below] at (3,-0.95) {P4};
  \draw[-{Latex[length=1.2mm]}, draw=primary-gold] (5,-0.95) -- (5,-0.5);
  \node[font=\tiny, text=primary-gold, below] at (5,-0.95) {P5};
\end{tikzpicture}
  \\[0.2cm]
  {\footnotesize\begin{tabular}{lrrrrrr}
\toprule
\textbf{Proc.} & \textbf{AT} & $\boldsymbol{\tau}$ & \textbf{CT} & \textbf{TAT} & \textbf{WT} & \textbf{RT} \\
\midrule
\texttt{P1} & 0 & 6 & 8 & 8 & 2 & 0 \\
\texttt{P2} & 1 & 8 & 25 & 24 & 16 & 16 \\
\texttt{P3} & 2 & 2 & 4 & 2 & 0 & 0 \\
\texttt{P4} & 3 & 5 & 17 & 14 & 9 & 9 \\
\texttt{P5} & 5 & 4 & 12 & 7 & 3 & 3 \\
\midrule
\multicolumn{4}{@{}l}{\textbf{Average}} & 11 & 6 & 5.6 \\
\bottomrule
\end{tabular}}
  \vspace{0.15cm}
  \begin{itemize}
    \footnotesize
    \item At 2, \texttt{P3} (2) beats \texttt{P1}'s remaining 4: preempt. At 3
      and 5, the newcomers (5 and 4) do not beat what is running.
      Average WT $= 30/5 = 6.0$~ms.
  \end{itemize}
\end{frame}

\begin{frame}{Three Policies, One Workload}
  \centering
  \footnotesize
  \renewcommand{\arraystretch}{1.3}
  \begin{tabular}{lcccc}
    \toprule
    \textbf{Policy} & \textbf{avg WT} & \textbf{avg TAT} & \textbf{avg RT} & \textbf{dispatches} \\
    \midrule
    FCFS & 9.2 & 14.2 & 9.2 & 5 \\
    SJF (non-preemptive) & 6.4 & 11.4 & 6.4 & 5 \\
    SRTF & 6.0 & 11.0 & 5.6 & 6 \\
    \bottomrule
  \end{tabular}
  \vspace{0.25cm}
  \begin{itemize}
    \item All three finish at 25~ms; the work is the same. Only the
      \emph{order} --- and hence the waiting --- changes.
    \item SRTF buys the best averages with one extra dispatch.
    \item With 0.5-ms dispatch latency, SRTF's utilization is
      $25/(25 + 6 \times 0.5) = 25/28 \approx 89.3\%$, versus $90.9\%$ for
      FCFS.
  \end{itemize}
  \muted{All times in ms; lab-server jobs from the opening slide.}
\end{frame}

\begin{frame}{The Price of Short-First: Starvation}
  \begin{itemize}
    \item If short jobs keep arriving, a long job can wait
      \emph{indefinitely}: SJF and SRTF can \key{starve} long processes.
    \item On the lab server, \texttt{P2} (8~ms) already has the largest WT
      (16~ms) under both SJF and SRTF; a steady stream of 2-ms jobs would
      postpone it forever.
    \item FCFS never starves anyone: each process waits only for those ahead
      of it.
  \end{itemize}
  \vspace{0.2cm}
  \begin{block}{Two further costs}
    SJF/SRTF need a burst prediction, and SRTF pays a dispatch on every
    preemption. Ageing (raising priority with waiting time) is the usual
    cure for starvation --- Lecture 9.
  \end{block}
\end{frame}

\begin{frame}{Common Misconceptions}
  \begin{columns}[T]
    \begin{column}{0.48\textwidth}
      \begin{block}{Wrong}
        \begin{itemize}
          \item ``Waiting time is completion minus arrival.''
          \item ``SJF uses the total length of the job.''
          \item ``SRTF preempts on equal remaining time.''
          \item ``A better policy finishes all the work sooner.''
        \end{itemize}
      \end{block}
    \end{column}
    \begin{column}{0.48\textwidth}
      \begin{block}{Right}
        \begin{itemize}
          \item That is TAT; $\mathrm{WT} = \mathrm{TAT} - \tau$.
          \item It uses the predicted \emph{next CPU burst}.
          \item Only a strictly shorter newcomer preempts.
          \item The last completion is the same; the averages change.
        \end{itemize}
      \end{block}
    \end{column}
  \end{columns}
\end{frame}

\begin{frame}{Takeaways}
  \begin{itemize}
    \item Processes alternate \key{CPU bursts} and \key{I/O bursts}; the
      short-term scheduler fills the CPU while others wait for I/O.
    \item Scheduling happens on block, exit (non-preemptive) and also on
      timer or I/O completion (\key{preemptive}); the \key{dispatcher} pays a
      latency on every switch.
    \item $\mathrm{TAT} = \mathrm{CT} - \mathrm{AT}$,
      $\mathrm{WT} = \mathrm{TAT} - \tau$,
      $\mathrm{RT} = \mathrm{ST} - \mathrm{AT}$.
    \item \key{FCFS} is simple and starvation-free but suffers the convoy
      effect.
    \item \key{SJF} minimizes average WT but needs a predicted burst
      (exponential average); \key{SRTF} is its preemptive form. Both can
      starve long jobs.
  \end{itemize}
\end{frame}

\begin{frame}{Next: Priority, Round Robin, and Multilevel Queues}
  \begin{itemize}
    \item Today: what scheduling is, how to measure it, and the two policies
      that order by arrival and by burst length.
    \item Next lecture: \key{Priority} scheduling and ageing, \key{Round
      Robin} with a time quantum, multilevel queues and feedback queues, and
      scheduling on more than one CPU.
  \end{itemize}
  \vspace{0.3cm}
  \begin{block}{Reading}
    Silberschatz et al., \emph{Operating System Concepts}, Ch.~5
    \S5.3--5.5 \cite{SilberschatzOSC_operating_system_concepts}; Stallings,
    \emph{Operating Systems: Internals and Design Principles}, Ch.~9
    \S9.2 and Ch.~10 \S10.1 \cite{StallingsOS_internals_design_principles}.
  \end{block}
\end{frame}

\begin{frame}{Further Reading (Page-Cited)}
  \footnotesize
  \begin{itemize}
    \item A.~Silberschatz, P.~B.~Galvin, G.~Gagne, \emph{Operating System
      Concepts}, 8e --- Ch.~5 \S5.1 (basic concepts, p.~183), \S5.2
      (scheduling criteria, p.~187), \S5.3.1 (FCFS, p.~188), \S5.3.2
      (SJF and SRTF, p.~189).
    \item W.~Stallings, \emph{Operating Systems: Internals and Design
      Principles}, 6e --- Ch.~9 \S9.1 (types of scheduling, p.~406), \S9.2
      (scheduling algorithms: FCFS, SPN, SRT, p.~410).
  \end{itemize}
  \vspace{0.2cm}
  \muted{Page-verified via each book's \texttt{index.md} under
  \texttt{master\_supporting\_docs/CS502/supporting\_books/}.}
\end{frame}

\begin{frame}[allowframebreaks]{References}
  \footnotesize
  \bibliographystyle{plain}
  \nocite{SilberschatzOSC_operating_system_concepts,StallingsOS_internals_design_principles}
  \bibliography{../../Bibliography_base}
\end{frame}

\end{document}
